\section{Diagnóstico do estágio atual}
\subsection{Estágio e evidências}
A PetCuida deve ser classificada como startup em \textbf{validação inicial}, anterior à tração. O repositório documenta protótipo funcional e telas de cadastro, histórico, lembretes, triagem orientativa e rede de parceiros. A Entrega XII e as versões anteriores reportam cinco sessões exploratórias guiadas; esse número não deve ser convertido em taxa de adoção nem em prova de retenção. Não há, nos materiais examinados, usuários pagantes, receita recorrente contratada ou clínicas com contrato comercial ativo.

O mercado brasileiro é relevante, mas o tamanho agregado do setor não equivale ao mercado endereçável da PetCuida. A Abempet reportou faturamento setorial de R\$ 77,96 bilhões em 2025 e projeção de aproximadamente R\$ 81 bilhões para 2026; esses números contextualizam o ambiente econômico, mas não demonstram demanda local por um aplicativo específico \cite{abempet2026,cnn2026}. A Pesquisa Nacional de Saúde de 2019 estimou cães em 33,8 milhões de domicílios e gatos em 14,1 milhões; também registrou vacinação antirrábica em 72\% dos domicílios com cães ou gatos, dado histórico que não deve ser apresentado como retrato atual \cite{ibge2020}.

\begin{table}[H]\centering
\caption{Diagnóstico do estágio da PetCuida}
\begin{tabularx}{\textwidth}{L{3.2cm}Y}
\toprule
Dimensão & Situação identificada nos documentos\\\midrule
Produto & Protótipo funcional; integrações externas, pagamentos, notificações e mapa real ainda exigem validação técnica.\\
Usuários & Cinco sessões exploratórias reportadas; amostra pequena e guiada, sem medida robusta de uso autônomo.\\
Mercado & O setor pet é expressivo, mas o segmento acessível e a demanda por praça não foram quantificados.\\
Receita & Preços e conversões são premissas; não há receita pagante comprovada.\\
Oferta & Clínicas e ONGs são potenciais parceiros; disponibilidade, adesão e disposição a pagar não estão contratadas.\\
Operação & Equipe de quatro integrantes, com capacidade limitada e competências concentradas.\\
\bottomrule
\end{tabularx}\fonteautores
\end{table}

\subsection{Desafios prioritários}
Os desafios que condicionam o crescimento são: (a) disposição a pagar ainda não observada; (b) CAC, ativação, retenção e conversão ainda não medidos; (c) dependência de convênio/patrocínio no cenário financeiro; (d) custo do fundo de créditos solidários; (e) problema de arranque bilateral entre tutores e clínicas; (f) segurança, privacidade e maturidade de integrações; e (g) capacidade operacional limitada. A ordem de prioridade deve ser reduzir incerteza sobre valor, oferta e economia unitária antes de investir em aquisição ampla.